\fichatitulo{32 · FUNDAMENTO}{NeurIPS · 2020}{Language Models are Few-Shot Learners}
{\small Brown et al. \textit{Language Models are Few-Shot Learners}. NeurIPS · 2020. \href{https://proceedings.neurips.cc/paper/2020/file/1457c0d6bfcb4967418bfb8ac142f64a-Paper.pdf}{Fonte consultada}.\par}

\campo{Problema e objetivo}
Até que ponto exemplos no contexto permitem executar tarefas sem ajuste de parâmetros?

\campo{Principal contribuição · síntese interpretativa}
Apresenta GPT-3 e avalia aprendizagem em contexto nas condições zero-shot, one-shot e few-shot.

\campo{Método / natureza da evidência}
Modelo autorregressivo de 175 bilhões de parâmetros; tarefas recebem instruções e exemplos no prompt, sem atualizações de gradiente na avaliação.

\campo{Resultado ou argumento principal}
O desempenho varia entre tarefas e com a quantidade de exemplos; o estudo demonstra a utilidade da adaptação por contexto.

\campo{Excertos-chave · seleção literal}
\begin{itemize}
\excerto{1}{without any gradient updates or fine-tuning}{Resumo.}
\excerto{2}{few-shot demonstrations}{Resumo.}
\end{itemize}

\campo{Limitações e análise crítica}
Análise da ficha: resultados em benchmarks não asseguram generalização ao domínio parlamentar. Exemplos, distribuição de teste e possível contaminação precisam ser controlados.

\campo{Aproveitamento na dissertação · proposta}
Fundamentação e metodologia: distinguir demonstrações no prompt de fine-tuning e registrar exemplos usados.

\registro{Fonte consultada nesta preparação: Resumo e introdução. Chave: \texttt{brownEtAl2020fewShotLearners}.}
